\chapter{الفضاءات التامة: بير، أسكولي، ستون--فايرشتراس}
\label{ch:b3:complete}

التمام --- أي تقارب كل متتالية كوشي --- هو الخاصية التي تتيح
للتحليل أن \emph{يُنتج} موضوعات: النقط الصامدة للتقلّصات، ومجاميع
المتسلسلات، وحلول المعادلات المحصَّل عليها بوصفها نهايات. ويجمع
هذا الفصل آلات الوجود الثلاث الكبرى في النظرية المترية.
إذ إن \emph{مبرهنة بير} تُظهر أن فضاءً \omterm{def:b3:complete:complete}{تامًّا} لا يمكن أن يكون اتحادًا
قابلًا للعدّ لقطع مهملة، وتستحضر موضوعات (دوال متصلة غير قابلة
للاشتقاق في أي نقطة!) من محض استدلال يشبه استدلالات عدد العناصر.
و\emph{أرزيلا--أسكولي} تعيّن الأجزاء المتراصة من $\mathcal C(K)$
وهي حصان التراص في التحليل --- وتستعملها مسألة نهاية الأسبوع
لإثبات مبرهنة بيانو للوجود في المعادلات التفاضلية.
و\emph{ستون--فايرشتراس} تُظهر أن كثيرات الحدود، وأشياء أخرى
كثيرة، كثيفةٌ في $\mathcal C(K)$: فيصير التقريب تحقّقًا جبريًا.
وفي الطريق نبني الإتمامات ونبرهن على مبرهنة التمديد للتطبيقات
المتصلة بانتظام، وهي الخبز اليومي في \cref{ch:b3:lp,ch:b3:hilbert,ch:b3:fouriertransform}.

\section{الفضاءات التامة والإتمامات والتمديدات}

\begin{definition}\label{def:b3:complete:complete}
يكون فضاء متري \emph{تامًّا}\index{فضاء متري تام} إذا تقاربت كل
متتالية كوشي (السنة الجامعية 2: الفضاء $\R^n$ تام؛ والفضاء
$\mathcal C(\intcc01)$ مع $\norm\cdot_\infty$ تام). وكل جزء مغلق من فضاء تام تامٌّ؛ وكل جزء تام
من أي فضاء متري مغلقٌ.
\end{definition}

\begin{proof}
من أجل العبارتين: تتقارب متتالية كوشي من المغلقة $F$ في $X$،
وتقع نهايتها، الملتصقة بالمجموعة $F$، في $F$؛ وتكون متتالية متقاربة في
$X$ من جزء تام $A$ متتاليةَ كوشي، فتتقارب في $A$، والنهايات
وحيدة.
\end{proof}

\begin{theorem}[تمديد التطبيقات المتصلة بانتظام]
\label{thm:b3:complete:extension}
لتكن $D \subseteq X$ كثيفة، وليكن $Y$ \omterm{def:b3:complete:complete}{تامًّا}، وليكن $f \colon D \to
Y$ \omterm{def:b3:topology:continuity}{متصلًا}
بانتظام. عندئذٍ يمتد $f$ \emph{بطريقة وحيدة} إلى \omterm{def:b3:topology:continuity}{تطبيق متصل}
$\bar f \colon X \to Y$، ويكون $\bar f$ \omterm{def:b3:topology:continuity}{متصلًا} بانتظام.\index{مبرهنة التمديد
(الاتصال المنتظم)}
\end{theorem}

\begin{proof}
الوحدانية: يتوافق تمديدان متصلان على $D$ الكثيفة، ومنه في كل
مكان (لأن مجموعة التوافق $\{g = h\}$ مغلقة: بوصفها سابق القطر
المغلق بالتطبيق $x\mapsto(g(x), h(x))$). الوجود: من أجل $x \in X$ نختار
$d_n \to x$ مع $d_n \in D$. وتكون المتتالية $(f(d_n))$ متتالية
كوشي: إذ إذا أُعطي $\varepsilon$، أعطى الاتصال المنتظم $\delta$
يحقق $d(u,v) < \delta \Rightarrow
d(f(u), f(v)) < \varepsilon$، والمتتالية $(d_n)$ متتالية كوشي. نضع $\bar f(x) = \lim f(d_n)$؛ ولا
تتعلق النهاية بالمتتالية المختارة (بتشبيك متتاليتين). ويمدّد
$\bar f$ التطبيقَ $f$ (بالمتتاليات الثابتة) ويرث معامل اتصاله:
إذ إذا كان $d(x, x') < \delta$، أعطى تقريب كليهما بنقط من $D$ على مسافة
$< \frac{\delta - d(x,x')}2$ أن $d(\bar f(x), \bar
f(x')) \leq \varepsilon$ بالنهاية --- أي إن $\bar f$ متصل بانتظام.
\end{proof}

\begin{theorem}[الإتمام]\label{thm:b3:complete:completion}
كل فضاء متري $X$ ينغمر تقييسيًا بوصفه جزءًا كثيفًا من \omterm{def:b3:complete:complete}{فضاء متري
تام} $\hat X$، وحيدٍ إلى حدود تقييس يُثبّت $X$ نقطةً نقطة: وهو
\emph{إتمامه}\index{إتمام فضاء متري}.
\end{theorem}

\begin{proof}
\emph{الوجود.} لتكن $\mathcal X$ مجموعة متتاليات كوشي في $X$،
مزوَّدةً بشبه المسافة
\[
D\bigl((x_n), (y_n)\bigr) = \lim_n d(x_n, y_n),
\]
والنهاية موجودة لأن $\abs{d(x_n, y_n) - d(x_m, y_m)}
\leq d(x_n, x_m) + d(y_n, y_m)$ يجعل المتتالية الحقيقية متتالية
كوشي. نضع $\hat X = \mathcal X/{\sim}$، بمطابقة المتتاليات التي المسافة $D$ بينها
معدومة؛ فينزل $D$ إلى مسافة. ونغمر $X$ بالمتتاليات الثابتة: وهو
تقييس، صورته كثيفة (إذ تُقرَّب متتالية كوشي بالمسافة $D$
بالثوابت المبنية على حدودها نفسها: $D\bigl((x_n), (x_k)_{\rm
const}\bigr) = \lim_n d(x_n, x_k) \to 0$ حين $k \to \infty$
بصفة كوشي). وأما تمام $\hat X$: فلتكن $(\xi^k)$ متتالية كوشي في
$\hat X$؛ نختار بالكثافة $x_k \in X$ يحقق $D(\xi^k, x_k)
\leq 2^{-k}$؛ عندئذٍ تكون
$(x_k)$ متتالية كوشي في $X$ (بمتراجحة المثلث عبر العناصر
$\xi$)، فتعرّف نقطة $\xi \in \hat
X$، ويكون $D(\xi^k, \xi) \leq 2^{-k} + D(x_k, \xi) \to 0$ (لأن المسافة من الثابت
$x_k$ إلى صنف $(x_j)_j$ هي $\lim_j d(x_k, x_j)$، وهي صغيرة من أجل $k$ كبير).

\emph{الوحدانية}: يحتوي إتمامان $\hat X_1, \hat X_2$ على $X$ كثيفةً؛ ويكون
تطابق $X$، وهو تقييس، \omterm{def:b3:topology:continuity}{متصلًا} بانتظام، ومنه يمتد إلى $\hat X_1 \to \hat X_2$
(\cref{thm:b3:complete:extension})، ويبقى تقييسًا على \omterm{def:b3:topology:interior}{مجموعة كثيفة} ومنه في كل مكان؛
وبالتناظر في الاتجاه الآخر، ويُثبّت المركّبان $X$ الكثيفة: أي
إنهما التطابقان.
\end{proof}

\begin{theorem}[النقطة الصامدة لباناخ]\label{thm:b3:complete:banach}
ليكن $X$ \omterm{def:b3:complete:complete}{تامًّا} وغير خالٍ وليكن $f \colon X \to X$ \emph{تقلّصًا}: أي $d(f(x), f(y)) \leq k\,d(x,y)$
مع $k < 1$. عندئذٍ للتطبيق $f$ نقطة صامدة وحيدة $x^*$، ويتقارب كل مدار
إليها، بالسرعة الصريحة $d(x_n, x^*) \leq
\frac{k^n}{1-k}\,d(x_1, x_0)$.\index{مبرهنة النقطة الصامدة
لباناخ}
\end{theorem}

\begin{proof}
(برهنت السنة الجامعية 2 على ذلك؛ ونعيد تسجيل الحجّة ذات السطرين
اكتفاءً ذاتيًا.) يحقق المدار $x_{n+1} = f(x_n)$ العلاقةَ $d(x_{n+1},
x_n) \leq k^nd(x_1, x_0)$، ومنه فهو
متتالية كوشي (بمتسلسلة هندسية)؛ ونهايته $x^*$ صامدة (باتصال
$f$)، ووحيدة لأن نقطتين صامدتين تحققان $d \leq k\,d$. وأما
السرعة: فبجمع الذيل الهندسي.
\end{proof}

\begin{example}[اضطراب التطابق]
\label{ex:b3:complete:perturbident}
ليكن $g \colon \R^d \to \R^d$ ليبشيتزيًا بثابت $k$ مع $k < 1$. عندئذٍ يكون $\varphi = \mathrm{id} + g$
\emph{تماثلًا طوبولوجيًا من $\R^d$ على $\R^d$}. والتباين، مع
معامل كمّي:
\[
\norm{\varphi(x) - \varphi(y)} \geq \norm{x - y} -
\norm{g(x) - g(y)} \geq (1 - k)\norm{x - y} .
\]
وأما الشمولية فهي مبرهنة النقطة الصامدة: إذ إن حل
$\varphi(x) = y$ يعني $x = y - g(x)$، ويكون $x \mapsto y -
g(x)$ تقلّصًا
بثابت $k$ للفضاء التام $\R^d$ --- ومنه يوجد حل وحيد
$x = \psi(y)$ من أجل كل $y$. وتجعل المتراجحة المعروضة المقلوبَ
$\psi$ ليبشيتزيًا بثابت $\frac1{1-k}$: أي تماثلًا طوبولوجيًا،
بحدود صريحة على المعاملين كليهما. وهذه العبارة التي تبدو
بريئة هي المحرّك داخل مبرهنة التطبيق العكسي (\cref{ch:b3:submanifolds}): إذ
يكون $f$، قرب نقطة يكون فيها $Df$ قابلًا للقلب، \emph{هو}
تطبيقًا خطيًا قابلًا للقلب مضافًا إليه اضطراب ليبشيتزي صغير،
ويقوم مثال اليوم بالباقي. وهي تُكمّم أيضًا المتانة العددية:
فجملة مضطربة بأقل من هامش المقلوب تبقى \omterm{def:b3:groups:derived}{قابلة للحل}، ويتحرّك
حلها بمقدار لا يتجاوز $\frac{1}{1-k}$ مضروبًا في الاضطراب.
\end{example}

\section{مبرهنة بير}

\begin{theorem}[بير]\label{thm:b3:complete:baire}
في \omterm{def:b3:complete:complete}{فضاء متري تام}، يكون كل تقاطع قابل للعدّ لمجموعات مفتوحة كثيفة
كثيفًا. وبصيغة مكافئة: إذا كان $X = \bigcup_{n}F_n$ مع كل $F_n$ مغلقة، فإن
داخل $F_n$ ما غير خالٍ.\index{مبرهنة بير}
\end{theorem}

\begin{proof}
لتكن $(U_n)$ مفتوحات كثيفة ولتكن $B_0 = B(x_0, r_0)$ أي كرة مفتوحة؛ سنجد
نقطة من $\bigcap U_n$ في $B_0$. وبالتتابع: بما أن $U_{n}$ كثيفة
ومفتوحة، فهي تلاقي الكرة المفتوحة $B_{n-1}$ في \omterm{def:b3:topology:topology}{مجموعة مفتوحة}،
وهذه تحتوي على كرة مغلقة $\bar B(x_n, r_n)$ مع $0 < r_n \leq r_{n-1}/2$ و $\bar B(x_n, r_n) \subseteq
B_{n-1}\cap U_n$. وتشكّل
المراكز متتالية كوشي (لأن $x_m \in
B_n$ من أجل $m \geq n$، وأنصاف
الأقطار $\to 0$)؛ وتقع النهاية $x$ في كل $\bar B(x_n, r_n)$ (بالانغلاق)،
ومنه في كل $U_n$ وفي $B_0$. وأما الصيغة الثانية: فإذا لم يكن
لأي $F_n$ داخل، كانت $U_n = X \setminus F_n$ مفتوحة وكثيفة، وأفلتت نقطة من
$\bigcap U_n$ من $\bigcup F_n = X$: وهو سخف.
\end{proof}

\begin{remark}\label{rem:b3:complete:meagre}
مفردات: نقول عن مجموعة إنها \emph{غير كثيفة في أي مكان} إذا كان
داخل غلقها خاليًا، وإنها \emph{هزيلة}\index{مجموعة هزيلة} (من
الصنف الأول) إذا كانت اتحادًا قابلًا للعدّ لمجموعات غير كثيفة في أي
مكان. وتقول مبرهنة بير: \emph{إن \omterm{def:b3:complete:complete}{الفضاء المتري التام} ليس هزيلًا
في نفسه}، وإن متمم \omterm{rem:b3:complete:meagre}{مجموعة هزيلة} كثيفٌ. و«الهزال» مفهوم للصغر
متعامد مع القياس (وسينتج \cref{ch:b3:measure} مجموعات \omterm{rem:b3:complete:meagre}{هزيلة} ذات قياس
كامل)، وتبرهن حجج بير على \emph{الوجود بالوفرة}: فلإظهار موضوع
واحد لا يحقق الخاصية P، بيّن أن الموضوعات المحققة للخاصية P تشكّل
\omterm{rem:b3:complete:meagre}{مجموعة هزيلة}.
\end{remark}

\begin{corollary}\label{cor:b3:complete:bairecors}
(a) المجموعة $\R$ غير قابلة للعدّ. (b) المجموعة $\Q$ ليست تقاطعًا
قابلًا للعدّ لأجزاء مفتوحة من $\R$، وكل \omterm{def:b3:complete:complete}{فضاء متري تام} غير خالٍ وبلا
نقط معزولة غيرُ قابل للعدّ.
\end{corollary}

\begin{proof}
(a) إن $\R = \bigcup_{x}\{x\}$ على مجموعة قابلة للعدّ سيجعل داخل مجموعة أحادية ما غير
خالٍ. (b) إذا كان $\Q = \bigcap_n V_n$ مع $V_n$ مفتوحة (وهي كثيفة بالضرورة،
لأنها $\supseteq \Q$)، فإن المجموعات $V_n$ والمتممات $\R\setminus\{q\}$ مع
$q \in \Q$ تشكّل عائلة قابلة للعدّ من المفتوحات الكثيفة تقاطعها
خالٍ --- وهذا يناقض بير. وإذا كان $X$ \omterm{def:b3:complete:complete}{تامًّا} بلا نقط معزولة
وقابلًا للعدّ، أظهره $X = \bigcup_{x \in X}\{x\}$ بوصفه اتحادًا قابلًا للعدّ لمجموعات مغلقة
داخلها خالٍ (لعدم وجود نقط معزولة): وهي مبرهنة بير مرة أخرى.
\end{proof}

\begin{theorem}[وحوش فايرشتراس موجودة]
\label{thm:b3:complete:nowherediff}
توجد دوال متصلة على $\intcc01$ غير قابلة للاشتقاق في أي نقطة.
بل إن مجموعة الدوال $f \in \mathcal C(\intcc01)$ التي لها مشتق (منتهٍ) ولو في نقطة
واحدة هزيلةٌ في $(\mathcal C(\intcc01), \norm\cdot_\infty)$.\index{دالة غير قابلة للاشتقاق في أي
نقطة}
\end{theorem}

\begin{proof}
من أجل $n \geq 1$ نضع
\[
F_n = \Bigl\{f : \exists x \in \intcc01,\
\forall h \neq 0 \text{ يحقق } x + h \in \intcc01,\
\abs{f(x+h) - f(x)} \leq n\abs h \Bigr\}.
\]
إذا كان $f$ قابلًا للاشتقاق عند $x$، فإن $f \in F_n$ من أجل $n$
ما: إذ إن النسبة $\abs{f(x+h)-f(x)}/\abs h$ محدودة من أجل $\abs h
\leq \delta$ (بالقابلية
للاشتقاق: فهي تؤول إلى $\abs{f'(x)}$) ومحدودة بالمقدار $2\norm f_\infty/\delta$ من أجل
$\abs h \geq \delta$. ومنه فإن $\bigcup F_n$ يحتوي على جميع الدوال القابلة
للاشتقاق في نقطة ما، ويكفي أن نبيّن أن كل $F_n$ مغلقة وداخلها
خالٍ.

\emph{مغلقة}: ليكن $f_k \to f$ بانتظام مع $f_k \in F_n$ بشواهد
$x_k \to x$ (بالتراص، بعد الاستخراج). ومن أجل $h$ يحقق $x + h \in \intcc01$:
نختار $h_k \to h$ مع $x_k + h_k
\in \intcc01$ (مثلًا $h_k = h + x - x_k$ مقصوصًا)؛ عندئذٍ
$\abs{f(x + h) - f(x)} = \lim \abs{f_k(x_k + h_k) - f_k(x_k)}
\leq \lim n\abs{h_k} = n\abs h$، باستعمال التقارب المنتظم واتصال $f$ عند النقط المعنية:
ومنه $f \in F_n$.

\emph{داخلها خالٍ}: إذا أُعطي $f \in F_n$ و $\varepsilon > 0$، وجدنا $g$
يحقق $\norm{g - f}_\infty \leq \varepsilon$ و $g
\notin F_n$. نقرّب أولًا $f$ في حدود $\varepsilon/2$
بدالة تآلفية بالقطع $\varphi$ (بالاتصال المنتظم: بالاستيفاء على
شبكة دقيقة)، ميولها محدودة بعدد $M$ ما. ثم نضيف منشارًا صغيرًا:
$g = \varphi + \frac\varepsilon2\,s_N$، حيث $s_N(x)$ الدالة المتعرّجة الدورية بالدور
$\frac1N$ وسعتها $1$ وميلها $\pm 2N$. وعند كل $x$، يوجد في أحد
الجانبين $h$ صغير كما نشاء يسهم فيه المنشار بميل $\pm 2N\cdot\frac\varepsilon2 =
\pm\varepsilon N$ على
$[x, x+h]$: ومنه تتجاوز نسبة تزايد $g$ المقدارَ $\varepsilon N - M > n$ من أجل
$N$ كبير. إذن $g \notin
F_n$، على أي مسافة منتظمة $\leq \varepsilon$ من $f$.
الخلاصة: المجموعة $\bigcup F_n$ \omterm{rem:b3:complete:meagre}{هزيلة}؛ ومنه فمتممها بمبرهنة بير
--- وهو مؤلَّف من الدوال غير القابلة للاشتقاق في أي نقطة ---
كثيفٌ في $\mathcal C(\intcc01)$: أي إن هذه الدوال موجودة \emph{بوفرة}.
\end{proof}

\section{أرزيلا--أسكولي}

في كل ما يلي، $K$ فضاء متري \emph{متراص} و $\mathcal
C(K) = \mathcal C(K, \R^d)$، مع $\norm f_\infty = \sup_K
\norm{f(x)}$:
وهو فضاء تام (لأن النهايات المنتظمة للدوال المتصلة متصلة ---
السنة الجامعية 2).

\begin{definition}\label{def:b3:complete:equicontinuous}
تكون عائلة $\mathcal F \subseteq \mathcal C(K)$ \emph{متساوية
الاتصال}\index{تساوي الاتصال} إذا وُجد، من أجل كل $\varepsilon > 0$،
عددٌ $\delta > 0$ بحيث
\[
d(x, y) < \delta \implies \norm{f(x) - f(y)} < \varepsilon
\quad\text{من أجل \emph{كل} } f \in \mathcal F
\]
(أي $\delta$ واحد للعائلة كلها --- مثل أي عائلة لها ثابت
ليبشيتز مشترك، أو معامل هولدر مشترك)، وتكون \emph{محدودة نقطةً
نقطة} إذا كان $\sup_{f}\norm{f(x)} < \infty$ من أجل كل $x$.
\end{definition}

\begin{theorem}[أرزيلا--أسكولي]\label{thm:b3:complete:ascoli}
يكون جزء $\mathcal F \subseteq \mathcal C(K)$ متراصًّا نسبيًا (أي غلقه متراص) إذا وفقط إذا
كان متساوي الاتصال ومحدودًا نقطةً نقطة. وعلى وجه الخصوص، لكل
\emph{متتالية} \omterm{def:b3:complete:equicontinuous}{متساوية الاتصال} ومحدودة نقطةً نقطة متتاليةٌ
جزئية متقاربة بانتظام.\index{مبرهنة أرزيلا--أسكولي}
\end{theorem}

\begin{proof}
($\Leftarrow$) لتكن $(f_k)$ متتالية من $\mathcal F$. الفضاء
المتري المتراص $K$ \emph{قابل للفصل}: إذ يغطّيه، من أجل كل $n$،
عددٌ منتهٍ من الكرات ذات نصف القطر $\frac1n$ (\cref{thm:b3:topology:metriccompact})؛
وتشكّل مراكزها مجموعة قابلة للعدّ كثيفة $D = \{x_1, x_2, \dots\}$. وبالحد نقطةً نقطة
وببولتزانو--فايرشتراس، نستخرج بالتتابع متتاليات جزئية تتقارب
عند $x_1$، ثم عند $x_2$ أيضًا، وهكذا، ونأخذ المتتالية الجزئية
\emph{القطرية} $(g_j)$: فهي تتقارب عند كل نقطة من $D$. ويرفع
\omterm{def:b3:complete:equicontinuous}{تساوي الاتصال} ذلك إلى صفة كوشي المنتظمة: إذ إذا أُعطي
$\varepsilon$، أخذنا $\delta$ كما في التعريف، وغطّينا $K$ بعدد
منتهٍ من الكرات $B(x_i, \delta)$ مع $x_i \in D$ (حيث
$i \leq m$)، واخترنا $J$ كبيرًا بحيث $\norm{g_j(x_i) - g_{j'}(x_i)} < \varepsilon$ من أجل $j, j' \geq
J$
و $i \leq m$. ومن أجل $x \in B(x_i, \delta)$ كيفي:
\[
\norm{g_j(x) - g_{j'}(x)}
\leq \norm{g_j(x) - g_j(x_i)} + \norm{g_j(x_i) - g_{j'}(x_i)}
+ \norm{g_{j'}(x_i) - g_{j'}(x)} < 3\varepsilon .
\]
إذن $(g_j)$ متتالية كوشي بانتظام، وتتقارب في الفضاء التام
$\mathcal C(K)$. ومن ثَمّ فإن لكل متتالية من $\mathcal F$
متتاليةً جزئية متقاربة: أي إن $\bar{\mathcal F}$ متراص (بالمتتاليات، ومنه
حسب \cref{thm:b3:topology:metriccompact}).

($\Rightarrow$) إذا كان $\bar{\mathcal F}$ متراصًّا: فالحد نقطةً نقطة واضح
(لأن التقييم متصل). وأما \omterm{def:b3:complete:equicontinuous}{تساوي الاتصال}، فنغطّي $\mathcal F$
بعدد منتهٍ من الكرات $B(f_i, \varepsilon)$ من $\mathcal C(K)$؛ وكل $f_i$ متصل
بانتظام (هاين، \cref{cor:b3:topology:heineborel})، فيعطي $\delta$ مشتركًا من أجل
$i \leq m$؛ وعندئذٍ من أجل $f \in B(f_i, \varepsilon)$ و $d(x,y)
< \delta$: $\norm{f(x)-f(y)} \leq 2\varepsilon +
\norm{f_i(x)-f_i(y)} < 3\varepsilon$.
\end{proof}

\begin{example}\label{ex:b3:complete:ascoliexamples}
الكرة الواحدية المغلقة من $\mathcal C(\intcc01)$ \emph{ليست} متراصة (فليس
للمتتالية $f_n(x) = x^n$ متتالية جزئية متقاربة بانتظام: لأن النهاية
نقطةً نقطة غير متصلة)، وبالفعل فإن $(x^n)$ ليست \omterm{def:b3:complete:equicontinuous}{متساوية الاتصال}
عند $1$. وفي المقابل، المجموعة $\{f : \norm
f_\infty \leq 1,\ \operatorname{Lip}(f) \leq 1\}$ متراصة: فهي محدودة
\omterm{def:b3:complete:equicontinuous}{ومتساوية الاتصال} بثابت ليبشيتز $1$، وهي مغلقة. وتفسّر مبرهنة
أسكولي \emph{لماذا} يخفق التراص في البُعد اللامنتهي (ريس، السنة
الجامعية 2) وما الذي يجب إضافته لاستعادته: معامل اتصال منتظم.
\end{example}

\section{ستون--فايرشتراس}

\begin{lemma}[ديني]\label{lem:b3:complete:dini}
ليكن $K$ متراصًّا ولتكن $(f_n)$ متتالية \emph{رتيبة} من الدوال
الحقيقية المتصلة تتقارب \emph{نقطةً نقطة} إلى دالة متصلة $f$.
عندئذٍ يكون التقارب منتظمًا.\index{مبرهنة ديني}
\end{lemma}

\begin{proof}
لنفترض $f_n \uparrow f$؛ ولتكن $g_n = f - f_n \downarrow 0$، وهي متصلة. إذا أُعطي
$\varepsilon$، فإن \omterm{def:b3:topology:topology}{المجموعات المفتوحة} $U_n = \{g_n <
\varepsilon\}$ متزايدة وتغطّي
$K$ (بالتقارب نقطةً نقطة)؛ نستخرج تغطية جزئية منتهية:
$K = U_{n_0}$ من أجل $n_0$ ما (لأن العائلة متزايدة)، أي $0 \leq g_n < \varepsilon$
في كل مكان من أجل $n \geq n_0$.
\end{proof}

\begin{lemma}\label{lem:b3:complete:sqrt}
توجد متتالية من \emph{كثيرات الحدود} $u_n$ تحقق $u_n(t) \to
\sqrt t$ بانتظام
على $\intcc01$.
\end{lemma}

\begin{proof}
نضع $u_0 = 0$ و $u_{n+1}(t) = u_n(t) + \frac12\bigl(t -
u_n(t)^2\bigr)$: وهي كثيرات حدود. وبالتراجع $0 \leq u_n(t) \leq
\sqrt t$ على
$\intcc01$: فبافتراض ذلك من أجل $n$،
\[
\sqrt t - u_{n+1}(t) = \bigl(\sqrt t - u_n(t)\bigr)
\Bigl(1 - \tfrac12\bigl(\sqrt t + u_n(t)\bigr)\Bigr) \geq 0,
\]
لأن $\sqrt t + u_n \leq 2$؛ و $u_{n+1} \geq u_n \geq 0$. إذن $(u_n(t))$ غير متناقصة ومحدودة
بالمقدار $\sqrt t$: فهي تتقارب نقطةً نقطة، وتحقق النهاية $\ell(t)$
العلاقةَ $\ell =
\ell + \frac12(t - \ell^2)$: أي $\ell(t) = \sqrt t$، وهي متصلة. وترفع مبرهنة ديني
(\cref{lem:b3:complete:dini}) ذلك إلى التقارب المنتظم.
\end{proof}

\begin{theorem}[ستون--فايرشتراس، الصيغة الحقيقية]
\label{thm:b3:complete:stoneweierstrass}
ليكن $K$ فضاءً متراصًّا (\omterm{def:b3:topology:hausdorff}{هاوسدورفيًا}) وليكن $\mathcal A \subseteq
\mathcal C(K, \R)$ \emph{جبرًا
جزئيًا} (صامدًا بالجمع والضرب والضرب السلَّمي) \emph{يحتوي على
الثوابت} و\emph{يفصل النقط} (أي إنه من أجل $x \neq y$ يوجد
$f \in
\mathcal A$ يحقق $f(x) \neq f(y)$). عندئذٍ يكون $\mathcal A$ كثيفًا
في $(\mathcal C(K,\R), \norm\cdot_\infty)$.\index{مبرهنة ستون--فايرشتراس}
\end{theorem}

\begin{proof}
ليكن $\bar{\mathcal A}$ الغلقَ، وهو جبر أيضًا (لأن جداءات النهايات
المنتظمة على مجموعات محدودة تتقارب).

\emph{الخطوة 1: $\bar{\mathcal A}$ شبكة}، أي صامدة بالعمليتين $\max$ و $\min$. وبما
أن $\max(f,g) = \frac{f + g +
\abs{f-g}}2$ وكذلك $\min$، يكفي أن يكون $f \in
\bar{\mathcal A} \Rightarrow \abs f \in \bar{\mathcal A}$: فبوضع $M = \norm f_\infty > 0$،
$\abs f = M\sqrt{(f/M)^2}$، تعطي \cref{lem:b3:complete:sqrt} كثيرات حدود $u_n$ تحقق $u_n\bigl((f/M)^2\bigr) \to \abs f/M$
بانتظام؛ وكثيرات الحدود في عناصر الجبر (مع الحد الثابت: لأن
الثوابت موجودة) تبقى في $\bar{\mathcal A}$.

\emph{الخطوة 2: الاستيفاء عند نقطتين.} من أجل $x \neq y$ ومن
أجل $a,
b \in \R$، يوجد $g \in \mathcal A$ يحقق $g(x) = a$ و $g(y) = b$: نأخذ
$h$ يفصل $x, y$ ونضع $g = a + (b -
a)\frac{h - h(x)}{h(y) - h(x)}$.

\emph{الخطوة 3.} ليكن $f \in \mathcal C(K)$ و $\varepsilon > 0$. من أجل كل زوج $x, y$
نختار $g_{x,y} \in \mathcal A$ يحقق $g_{x,y}(x) = f(x)$ و $g_{x,y}(y) = f(y)$ (الخطوة 2؛ ومن أجل $x = y$
نأخذ الدالة الثابتة $g_{x,x} = f(x)$). نثبّت $x$: فمن أجل كل
$y$، تحتوي \omterm{def:b3:topology:topology}{المجموعة المفتوحة} $V_y = \{g_{x,y} < f +
\varepsilon\}$ على $y$؛ ويستخرج التراص
$y_1, \dots,
y_m$ يحقق $K = \bigcup V_{y_j}$، ويحقق $h_x = \min_j g_{x, y_j}
\in \bar{\mathcal A}$ (الخطوة 1) العلاقةَ $h_x < f +
\varepsilon$
في كل مكان مع $h_x(x) = f(x)$. والآن نغيّر $x$: تكون $W_x =
\{h_x > f - \varepsilon\}$
مفتوحة وتحتوي على $x$؛ نستخرج $x_1,
\dots, x_l$ تغطّي $K$، وتحقق $h = \max_i h_{x_i} \in
\bar{\mathcal A}$
العلاقةَ $f - \varepsilon < h < f +
\varepsilon$: أي $\norm{f - h}_\infty \leq \varepsilon$. ومنه $f \in \bar{\mathcal A}$.
\end{proof}

\begin{corollary}\label{cor:b3:complete:weierstrass}
(a) (\emph{فايرشتراس}) كثيرات الحدود كثيفة في $\mathcal
C(\intcc ab, \R)$؛ وكثيرات
الحدود ذات $n$ متغيّرًا كثيفة في $\mathcal C(K, \R)$ من أجل $K \subseteq \R^n$ متراصة.
(b) (\emph{الصيغة العقدية}) إذا كان $\mathcal A \subseteq \mathcal
C(K, \C)$ جبرًا جزئيًا يحتوي
على الثوابت ويفصل النقط و\emph{صامدًا بالاقتران}، فهو كثيف.
(c) (\emph{الصيغة المثلثية}) كثيرات الحدود المثلثية $\sum_{\abs n \leq N}c_n\eu^{\iu n t}$
كثيفة في فضاء الدوال المتصلة الدورية بالدور $2\pi$ مزوَّدًا
بالمعيار $\norm\cdot_\infty$.\index{مبرهنة فايرشتراس للتقريب}
\end{corollary}

\begin{proof}
(a) كثيرات الحدود تشكّل جبرًا يحتوي على الثوابت؛ ودوال
الإحداثيات تفصل نقط $\R^n$.
(b) يشكّل الجزآن الحقيقي والتخيلي $\frac{f + \bar f}2$ و $\frac{f - \bar f}{2\iu}$ لعناصر
$\mathcal A$ جبرًا حقيقيًا $\mathcal A_\R \subseteq \mathcal C(K, \R)$ يحتوي على الثوابت؛ وهو يفصل
النقط (إذ يفرض $f(x) \ne f(y)$ أن يفصل $\operatorname{Re}f$ أو $\operatorname{Im}f$).
نطبّق المبرهنة الحقيقية ونعيد التركيب.
(c) ننظر إلى الدوال الدورية بالدور $2\pi$ على أنها $\mathcal C(S^1, \C)$
($S^1 = \R/2\pi\Z$، وهي متراصة: \cref{exo:b3:topology:5})؛ والجبر المولَّد بالعناصر $\eu^{\iu t}, \eu^{-\iu t}$
وبالثوابت صامد بالاقتران ويفصل نقط الدائرة (لأن $\eu^{\iu t}$
متباين عليها). نطبّق (b).
\end{proof}

\begin{remark}\label{rem:b3:complete:swremark}
تُصلح الصيغة المثلثية الثغرةَ التي تركها فصل فورييه في السنة
الجامعية 2 وتعمّمها تعميمًا واسعًا: فكثافة كثيرات الحدود
المثلثية في $(\mathcal C(S^1), \norm\cdot_2)$ تنتج من باب أولى ($\norm\cdot_2 \leq \norm\cdot_\infty$ إلى حدود ثابت
التنظيم)، وهذا ما سيجعل جملة فورييه أساسًا متعامدًا
\emph{ممنظّمًا} في \cref{ch:b3:hilbert}، فيبرهن على مبرهنة بارسيفال في
عمومها التام أخيرًا.
\end{remark}

\section{تمارين}

\begin{exercise}[$\star$]\label{exo:b3:complete:1}
(a) برهن على أن $\mathcal C(\intcc01, \R)$ مزوَّدًا بالمعيار $\norm f_1 = \int_0^1\abs f$ \emph{ليس} \omterm{def:b3:complete:complete}{تامًّا}:
فالدوال $f_n$ المساوية $0$ على $\intcc0{\frac12}$ و $1$ على $\intcc{\frac12 + \frac1n}1$
والتآلفية بينهما متتاليةُ كوشي بالمعيار $\norm\cdot_1$ وليست لها
نهاية متصلة.
(b) برهن على أن كل فضاء معياري تتقارب فيه كل متسلسلة متقاربة
مطلقًا هو فضاء تام. \emph{(استخرج من متتالية كوشي متتاليةً
جزئية تحقق $\norm{x_{n_{k+1}} -
x_{n_k}} \leq 2^{-k}$.)}
\end{exercise}

\begin{exercise}[$\star$]\label{exo:b3:complete:2}
باستعمال مبرهنة بير: (a) برهن على أنه ليس لفضاء معياري تام أساس
(\omterm{def:b3:galois:algebraic}{جبري}) قابل للعدّ --- واستنتج أن فضاء كثيرات الحدود ليس \omterm{def:b3:complete:complete}{تامًّا} من
أجل \emph{أي} معيار؛ (b) برهن على أنه إذا تقاربت متتالية من
الدوال المتصلة $f_n \colon \R \to \R$ نقطةً نقطة إلى $f$، فإن مجموعة نقط
اتصال $f$ كثيفة. \emph{(من أجل (b): اقبل أو برهن على أن $\Omega_\delta = \{x :
\operatorname{osc}_x f < \delta\}$
مفتوحة، وبرهن على أنها كثيفة باستعمال $F_{N} = \{x: \abs{f_n(x) - f_m(x)} \leq \delta/3\
\forall n,m \geq N\}$، وهي مجموعات
مغلقة تغطّي $\R$؛ واعمل في كرة مغلقة كيفية لتطبيق بير.)}
\end{exercise}

\begin{exercise}[$\star\star$]\label{exo:b3:complete:3}
(a) برهن على أن $f(x) = \frac12\bigl(x + \frac ax\bigr)$ (حيث $a >
1$) تقلّصٌ للمجموعة $[\sqrt a, +\infty)$،
وعيّن نقطته الصامدة --- وهي طريقة هيرون. وقدّر عدد التكرارات
اللازمة لدقة $10^{-12}$ انطلاقًا من $x_0 = a$، من أجل $a = 2$.
(b) (معادلة كبلر) من أجل $0 \leq e < 1$ و $m \in \R$، برهن على
أن $x = m + e\sin x$ لها حل وحيد يتعلق اتصالًا بالوسيط $m$.
\end{exercise}

\begin{exercise}[$\star\star$]\label{exo:b3:complete:4}
(a) ليكن $X$ \omterm{def:b3:complete:complete}{تامًّا} وليكن $f \colon X \to X$ بحيث تكون قوة ما $f^p$
تقلّصًا. برهن على أن للتطبيق $f$ نقطة صامدة وحيدة. تطبيق: يحقق المؤثر
التكاملي $T$ على $\mathcal C(\intcc0a)$، $Tf(x) = \int_0^x f$، العلاقةَ $\norm{T^p} \leq a^p/p!$ --- واستنتج
أن $u = g + \lambda Tu$ \omterm{def:b3:groups:derived}{قابلة للحل} من أجل كل $\lambda$.
(b) (إيدلشتاين) ليكن $K$ \emph{متراصًّا} وليكن $f\colon K \to K$ يحقق
$d(f(x), f(y)) < d(x, y)$ من أجل $x \neq y$. برهن على أن للتطبيق $f$ نقطة صامدة وحيدة،
لكن معدل التقلّص يمكن أن يضيع: فعلى $X = [1, +\infty)$ (وهو تام وغير
متراص)، ليس للتطبيق $f(x) = x
+ \frac1x$ نقطة صامدة رغم أنه ينقص المسافات نقصًا
\omterm{def:b3:complete:complete}{تامًّا}.
\end{exercise}

\begin{exercise}[$\star\star$]\label{exo:b3:complete:5}
(a) تطبيقان متصلان إلى فضاء هاوسدورفي يتوافقان على جزء كثيف
يتوافقان في كل مكان؛ وأين استُعمل ذلك في الفصل؟
(b) لتكن $D \subseteq X$ كثيفة وليكن $f \colon D \to Y$ تقابلًا تقييسيًا
على جزء كثيف من فضاء تام $Y$، مع $X$ تام. برهن على أن $f$ يمتد
إلى تقابل تقييسي $X \to Y$. واستنتج وحدانية الإتمامات من جديد.
\end{exercise}

\begin{exercise}[$\star\star$]\label{exo:b3:complete:6}
أي العائلات التالية \omterm{def:b3:complete:equicontinuous}{متساوية الاتصال}، ومحدودة نقطةً نقطة،
ومتراصة نسبيًا في $\mathcal C(\intcc01)$؟
\[
\{x \mapsto \sin(nx)\}_{n},\qquad
\{x \mapsto x^n\}_n,
\]
\[
\Bigl\{f \in \mathcal C^1 : \norm f_\infty \leq 1,\
\norm{f'}_\infty \leq 1\Bigr\},\qquad
\{f : \operatorname{Lip}(f)\leq 1,\ f(0) = 0\}.
\]
برّر كل جواب بمبرهنة أسكولي أو بمتتالية مضادّة.
\end{exercise}

\begin{exercise}[$\star\star\star$]\label{exo:b3:complete:7}
(المؤثرات التكاملية المتراصة) ليكن $k \in \mathcal C(\intcc01^2)$ ولنضع، من أجل
$f \in \mathcal C(\intcc01)$، $Tf(x) =
\int_0^1k(x,y)f(y)\,\dd y$.
(a) برهن على أن $T$ يرسل الكرة الواحدية من $\mathcal C(\intcc01)$ إلى مجموعة
\omterm{def:b3:complete:equicontinuous}{متساوية الاتصال} ومحدودة بانتظام؛ واخلص إلى أن $T$ \emph{مؤثر
متراص}: أي إن صور المجموعات المحدودة متراصة نسبيًا.
(b) استنتج أنه إذا كانت $(f_n)$ محدودة، فإن للمتتالية $(Tf_n)$ متتالية
جزئية متقاربة بانتظام، وأن $T$ لا يمكن أن يكون تقابلًا مقلوبه
متصل. \emph{(إذ ستكون صورة الكرة الواحدية جوارًا متراصًّا
للنقطة $0$ في $\mathcal C(\intcc01)$: وهو ممنوع بمبرهنة ريس من السنة الجامعية
2.)}
\end{exercise}

\begin{exercise}[$\star\star$]\label{exo:b3:complete:8}
برهن أو انقض، من أجل $f_n \colon \intcc01 \to \R$ متصلة:
(a) $f_n \downarrow 0$ نقطةً نقطة $\Rightarrow$ بانتظام (ديني --- أعد
البرهان عليها)؛ (b) الأمر نفسه دون الرتابة؛ (c) الأمر نفسه مع
الرتابة لكن مع $f$ غير متصلة؛ (d) الأمر نفسه مع الرتابة ونهاية
متصلة، لكن على $\intoo01$.
\end{exercise}

\begin{exercise}[$\star\star$]\label{exo:b3:complete:9}
(a) (العزوم تحدّد) لتكن $f \in \mathcal C(\intcc01)$ مع $\int_0^1 f(x)\,x^n\,\dd x = 0$ من أجل كل
$n \in \N$. برهن على $f =
0$. \emph{(قرّب $f$ بانتظام بكثيرات
حدود واحسب $\int f^2$.)}
(b) برهن على أن كثيرات الحدود الزوجية كثيفة في $\mathcal
C(\intcc01)$ لكنها
\emph{ليست} كثيفة في $\mathcal C(\intcc{-1}1)$؛ وأين تخفق فرضية
ستون--فايرشتراس؟
(c) هل الجبر المولَّد بالعنصر $x \mapsto \eu^{\iu x}$ وحده (دون $\eu^{-\iu x}$) كثيف
في $\mathcal C(S^1, \C)$؟ \emph{(انظر في $\int_0^{2\pi}f(t)\,\eu^{\iu t}\dd t$.)}
\end{exercise}

\begin{exercise}[$\star\star\star$]\label{exo:b3:complete:10}
(الحد المنتظم، الصيغة المترية) ليكن $X$ فضاءً متريًا \omterm{def:b3:complete:complete}{تامًّا}
ولتكن $\mathcal F \subseteq \mathcal C(X, \R)$ عائلة محدودة \emph{نقطةً نقطة}: أي $\sup_{f \in \mathcal
F}\abs{f(x)} < \infty$ من أجل
كل $x$. برهن على وجود مفتوحة غير خالية $U \subseteq X$ تكون
$\mathcal F$ محدودة عليها \emph{بانتظام}: أي $\sup_{f}\sup_U \abs f < \infty$.
\emph{(انظر في $F_n = \{x : \abs{f(x)} \leq n\ \forall f\}$.)} وهذا هو المحرّك وراء مبرهنة
باناخ--شتاينهاوس في \cref{ch:b3:banach}.
\end{exercise}

\begin{exercise}[$\star\star$]\label{exo:b3:complete:11}
(الصنف $\mathcal C^1$ يحتاج معياره الخاص) على $E = \mathcal
C^1(\intcc01, \R)$ لننظر في
$\norm f_{\mathcal C^1} = \norm
f_\infty + \norm{f'}_\infty$.
(a) برهن على أن $(E, \norm\cdot_{\mathcal C^1})$ تام \emph{(فلمتتالية كوشي بالمعنى
$\mathcal C^1$ يكون $f_n \to f$ و $f_n' \to g$ بانتظام؛ ونعيّن
$g = f'$ بالانتقال إلى النهاية في $f_n(x) = f_n(0) + \int_0^xf_n'$)}.
(b) برهن على أن $(E, \norm\cdot_\infty)$ \emph{ليس} \omterm{def:b3:complete:complete}{تامًّا}: أظهر نهاية منتظمة
لدوال من الصنف $\mathcal C^1$ ليست قابلة للاشتقاق (مثل تقريبات
ملساء للدالة $\abs{x - \tfrac12}$).
(c) استنتج من (a) و (b) ومن دائرة أفكار التطبيق المفتوح --- أو
مباشرةً --- أنه لا يوجد ثابت $C$ يحقق $\norm{f'}_\infty \leq C\,\norm f_\infty$ على $E$: أظهر
متتالية تشهد بذلك. فالاشتقاق غير محدود؛ وهذه هي الهاوية وراء
\cref{thm:b3:complete:nowherediff}.
\end{exercise}

\begin{exercise}[$\star\star\star$]\label{exo:b3:complete:12}
(مبرهنة كروفت المساعدة) لتكن $f\colon\intoo0\infty\to\R$ متصلة ولنفترض أنه من أجل
كل $x > 0$ يكون $f(nx) \to 0$ حين $n \to \infty$ عبر الأعداد
الصحيحة. برهن على أن $f(t) \to 0$ حين $t \to
+\infty$. \emph{(ثبّت
$\varepsilon > 0$؛ فالمجموعات
\[
F_N = \{x > 0 : \abs{f(nx)} \leq \varepsilon\ \ \forall n
\geq N\}
\]
مغلقة وتغطّي $\intoo0\infty$؛ وتعطي مبرهنة بير في فترة ما
$\intcc ab$ عددًا $N$ وفترةً جزئيةً $\intcc{a'}{b'}
\subseteq F_N$؛ وعندئذٍ تغطّي
المضاعفات $\bigl[na', nb'\bigr]$، $n
\geq N$، جوارًا كاملًا للنقطة $+\infty$ متى
كان $n(b' - a') \geq a'$.)} وأين استُعملت فرضية «من أجل كل $x$» (لا من أجل
$x$ الناطقة فحسب)؟
\end{exercise}

\section{مسألة: مبرهنة بيانو للوجود}

\begin{problem}[{مسألة نهاية الأسبوع --- وجود حلول المعادلة
$x' = f(t, x)$ دون شرط ليبشيتز}]\label{pb:b3:complete:1}
تتطلب مبرهنة كوشي--ليبشيتز (السنة الجامعية 2؛ ويُعاد البرهان
عليها في \cref{ch:b3:ode}) أن يكون $f$ ليبشيتزيًا في $x$. أما بيانو
(1890) فقد أثبت أن \emph{اتصال $f$ يعطي الوجود أصلًا} --- وإن
لم يعطِ الوحدانية. وسنبرهن على ذلك بمضلعات أويلر ومبرهنة
أسكولي. المعطيات: $f \colon R
\to \R^d$ متصل على المستطيل $R = \intcc{t_0 - a}{t_0
+ a}\times \bar B(x_0, b) \subseteq \R\times\R^d$، $M =
\sup_R\norm f$،
و
\[
T = \min\bigl(a,\ b/M\bigr)
\qquad (\text{إذا كان } M = 0 \text{ فالمسألة تافهة}).
\]

\medskip
\noindent\textbf{الجزء الأول --- مضلعات أويلر.} من أجل
$n \geq 1$، نجزّئ $[t_0, t_0 + T]$ بالنقط $t_k = t_0 + kT/n$ (حيث $0 \leq k
\leq n$)
ونعرّف $\varphi_n$ تآلفيًا بالقطع: $\varphi_n(t_0) = x_0$، وعلى
$[t_k, t_{k+1}]$،
\[
\varphi_n(t) = \varphi_n(t_k) + (t -
t_k)\,f\bigl(t_k, \varphi_n(t_k)\bigr).
\]
\begin{enumerate}
  \item برهن بالتراجع على أن $\varphi_n$ معرَّف تعريفًا سليمًا،
        مع $\norm{\varphi_n(t) - x_0} \leq M(t - t_0) \leq b$ على $[t_0, t_0 + T]$ --- ومنه تبقى نقط
        التقييم في $R$. \emph{(وهنا يدخل الشرط $T \leq b/M$.)}
  \item برهن على أن كل $\varphi_n$ ليبشيتزي بثابت $M$.
  \item استنتج من مبرهنة أرزيلا--أسكولي (\cref{thm:b3:complete:ascoli}) أن
        متتالية جزئية $(\varphi_{n_j})$ تتقارب بانتظام على $[t_0, t_0 +
        T]$ إلى
        دالة $\varphi$، وهي نفسها ليبشيتزية بثابت $M$ وتحقق
        $\varphi(t_0) = x_0$.
\end{enumerate}

\medskip
\noindent\textbf{الجزء الثاني --- النهاية تحل المعادلة.} نعرّف
\emph{القصور} $\Delta_n(t) = \varphi_n'(t) -
f\bigl(t, \varphi_n(t)\bigr)$ عند النقط خارج الشبكة.
\begin{enumerate}[resume]
  \item برهن على أن $f$ متصل بانتظام على $R$، واستنتج: أنه من
        أجل كل $\varepsilon > 0$ يوجد $n_0$ بحيث يكون، من أجل $n \geq n_0$
        ومن أجل كل $t$ خارج الشبكة، $\norm{\Delta_n(t)} \leq \varepsilon$. \emph{(فعلى $(t_k,
        t_{k+1})$،
        $\varphi_n'(t) = f(t_k, \varphi_n(t_k))$، ويقع $(t, \varphi_n(t))$ على مسافة لا تتجاوز
        $(1 + M)\,T/n$ من $(t_k, \varphi_n(t_k))$.)}
  \item أثبت الصيغة التكاملية: من أجل كل $t$،
        \[
        \varphi_n(t) = x_0 + \int_{t_0}^{t}
        f\bigl(s, \varphi_n(s)\bigr)\dd s +
        \int_{t_0}^t \Delta_n(s)\,\dd s,
        \]
        والمقدار المكامَل الأوسط متصل بالقطع.
  \item انتقل إلى النهاية على $(n_j)$: برهن على $f(s, \varphi_{n_j}(s)) \to f(s, \varphi(s))$
        \emph{بانتظام} (بالاتصال المنتظم للتطبيق $f$ مرة أخرى)،
        واخلص إلى
        \[
        \varphi(t) = x_0 + \int_{t_0}^t
        f\bigl(s, \varphi(s)\bigr)\dd s .
        \]
  \item استنتج أن $\varphi$ من الصنف $\mathcal C^1$ على $[t_0, t_0 +
        T]$
        ويحل $x' = f(t, x)$ مع $x(t_0) = x_0$؛ ومدّد البناء إلى
        $[t_0 - T, t_0]$ (بعكس الزمن). \emph{وهذه هي مبرهنة
        بيانو.}
\end{enumerate}

\medskip
\noindent\textbf{الجزء الثالث --- الوحدانية تخفق فعلًا.} لننظر
في $x' = 2\sqrt{\abs x}$ مع $x(0) = 0$ على $\R$.
\begin{enumerate}[resume]
  \item تحقق من أن $f(x) = 2\sqrt{\abs x}$ متصل لكنه ليس ليبشيتزيًا على أي \omterm{def:b3:topology:topology}{جوار}
        للنقطة $0$.
  \item تحقق من أن $x \equiv 0$، ومن أجل كل $c \geq 0$،
        \[
        x_c(t) = \begin{cases} 0 & t \leq c,\\ (t - c)^2 & t
        \geq c,\end{cases}
        \]
        كلها حلول مارّة بالنقطة $(0,0)$: أي متصلٌ من الحلول
        المتمايزة.
  \item وأين تنهار حجّة تكرار بيكار (النقطة الصامدة لباناخ) من
        أجل هذا $f$؟
\end{enumerate}

\medskip
\noindent\textbf{الجزء الرابع --- حدود الطريقة.}
\begin{enumerate}[resume]
  \item برهن على أن مبرهنة بيانو تخفق في البُعد اللامنتهي:
        نقبل (أو يمكنك أن تأخذه تسليمًا) مثالَ ديودونيه
        الكلاسيكي في فضاء المتتاليات المتلاشية $c_0$؛ وبدلًا من
        ذلك، برهن على المكوّن المنتهي البُعد الذي يخفق هناك:
        الكرة الواحدية المغلقة من $c_0$ (بمعيار السوپريموم) ليست
        متراصة --- أظهر متتالية محدودة بلا متتالية جزئية
        متقاربة، واشرح أي خطوة من الجزء الأول تنهار.
  \item لخّص: أي الفرضيات تعطي الوجود؟ وأيها يعطي الوجود
        والوحدانية؟ اذكر بدقة المبرهنتين (بيانو؛
        كوشي--ليبشيتز) جنبًا إلى جنب.
\end{enumerate}

\medskip
\noindent\textbf{الجزء الخامس --- أوسغود: الوحدانية أبعد من
ليبشيتز.} لتكن $\omega\colon\intoo0\infty\to\intoo0\infty$ متصلة وغير متناقصة وتحقق
\[
\int_0^1\frac{\dd r}{\omega(r)} = +\infty
\qquad(\text{تباعد عند } 0),
\]
ولنفترض أن $f$ يحقق $\norm{f(t, x) - f(t, y)} \leq
\omega\bigl(\norm{x - y}\bigr)$ على $R$.
\begin{enumerate}[resume]
  \item تحقق من أن $\omega(r) = Lr$ صالحة (ليبشيتز)، ومن أن
        $\omega(r) = r\log\frac1r$ (ممدَّدة بالاتصال، من أجل $r$ صغيرة) صالحة رغم
        أنها \emph{ليست} $O(r)$، ومن أن $\omega(r) = 2\sqrt r$ ليست كذلك. واحسب
        التكامل في كل حالة.
  \item ليكن $x_1, x_2$ حلين للمعادلة على $[t_0, t_0 + T]$ مع
        $x_1(t_0) = x_2(t_0)$، ولتكن $\delta(t) =
        \norm{x_1(t) - x_2(t)}$. برهن، انطلاقًا من الصيغ
        التكاملية وحدها، على أنه من أجل $t_0 \leq s \leq t$:
        \[
        \delta(t) \leq \delta(s) +
        \int_s^t\omega\bigl(\delta(v)\bigr)\dd v .
        \]
  \item (مبرهنة أوسغود) لنفترض $\delta(t_1) > 0$ من أجل $t_1$ ما، ولتكن
        $\tau = \sup\{t \leq t_1 : \delta(t) =
        0\}$. من أجل $s \in \intoo\tau{t_1}$، نضع $u(t) = \delta(s)
        + \int_s^t\omega(\delta(v))\,\dd v$. برهن على $\delta \leq
        u$
        و $u' = \omega(\delta) \leq \omega(u)$، واستنتج
        \[
        \int_{\delta(s)}^{u(t_1)}\frac{\dd r}{\omega(r)}
        \leq t_1 - s \leq t_1 - \tau .
        \]
        ثم اجعل $s \downarrow \tau$ واستخرج تناقضًا مع تباعد التكامل. واخلص
        إلى: \emph{أن الحلول المارّة بشرط ابتدائي مشترك
        تتطابق} --- أي الوحدانية تحت شرط أوسغود.
  \item استخلص النتائج: فوحدانية كوشي--ليبشيتز هي الحالة
        $\omega(r) = Lr$؛ وللمعادلة $x' =
        x\log\frac1{\abs x}$ (ممدَّدة بالقيمة $0$
        عند $0$) حلولٌ وحيدة رغم أن طرفها الأيمن ليس ليبشيتزيًا
        عند $0$؛ ومن أجل $x' = 2\sqrt{\abs x}$ يكون تقارب $\int_0\frac{\dd r}{2\sqrt r}$ هو بالضبط
        ما يتيح لحل أن يغادر $0$ في زمن منتهٍ --- طابق قيمة
        التكامل $\int_0^{h^2}\frac{\dd r}{2\sqrt r} = h$ مع سلوك الهروب للحل $x_c$.
\end{enumerate}

\medskip
\noindent\textbf{الجزء السادس --- السرعات والمخططات والأقماع.}
\begin{enumerate}[resume]
  \item (مبرهنة غرونوال التكاملية المساعدة) ليكن $e, \eta \geq 0$
        متصلين على $[t_0, t_0+T]$ مع $e(t) \leq \int_{t_0}
        ^t\bigl(L\,e(s) + \eta(s)\bigr)\dd s$ من أجل كل $t$. برهن
        على
        \[
        e(t) \leq \frac{\sup\eta}{L}\bigl(\eu^{L(t - t_0)} -
        1\bigr)
        \]
        \emph{(ضع $G(t) = \int_{t_0}^t(Le + \eta)$، ولاحظ $G'
        \leq LG + \sup\eta$، واشتق $\eu^{-Lt}G(t)$)}.
  \item (أويلر يتقارب بسرعة) لنفترض الآن أن $f$ ليبشيتزي بثابت
        $L$ في $x$ وبثابت $L'$ في $t$ على $R$. وبالجمع بين حدّ
        القصور في السؤال 4 (مجعولًا كمّيًا: $\norm{\Delta_n} \leq (L' + LM)T/n$) والسؤال 17،
        برهن على
        \[
        \sup_{[t_0, t_0+T]}\norm{\varphi_n - \varphi}
        \;\leq\; \frac{(L' + LM)\,T}{n}\cdot
        \frac{\eu^{LT} - 1}{L} = O\Bigl(\frac1n\Bigr),
        \]
        حيث $\varphi$ هو \emph{الحل} الوحيد: أي إن المتتالية
        كلها تتقارب مع معطيات ليبشيتزية، بسرعة صريحة --- دون
        حاجة إلى متتاليات جزئية. ولماذا ترفع الوحدانية التقاربَ
        الجزئي إلى تقارب كامل حتى دون هذا الحساب؟
  \item (المخطط يختار) من أجل $x' = 2\sqrt{\abs x}$، $x(0) =
        0$: برهن على أن كل
        مضلع أويلر معدوم تطابقيًا، ومنه يتقارب المخطط إلى الحل
        $x \equiv 0$؛ لكنه، إذا انطلق من $x(0) = \varepsilon > 0$، يتقارب (حين
        $n \to \infty$ ثم $\varepsilon \to 0$) إلى $t
        \mapsto t^2$، وهو حل
        \emph{مختلف} مارّ بالمبدأ. فتظهر عدم الوحدانية من جديد
        على صورة حساسية المخطط العددي للاضطرابات.
  \item برهن على أن مجموعة $\mathcal S$ المؤلَّفة من \emph{جميع}
        حلول $x' = f(t,x)$ مع $x(t_0) = x_0$ على $[t_0, t_0+T]$
        (بقيم في $\bar B(x_0, b)$) غير خالية (الجزء الثاني)،
        وليبشيتزية بانتظام بثابت $M$، ومغلقة في $\bigl(\mathcal C([t_0, t_0+T], \R^d),
        \norm\cdot_\infty\bigr)$؛ واخلص
        بمبرهنة أسكولي إلى أن $\mathcal S$ \emph{متراصة}. (وتضيف
        مبرهنة كنيزر أن $\mathcal S$ مترابطة؛ ولن نبرهن عليها.)
  \item تحقق من ظاهرة كنيزر على المثال: من أجل $x' =
        2\sqrt{\abs x}$ مع
        $x(0) = 0$ على $[0, 1]$، برهن على $\mathcal S = \{x_c : c \in \intcc0\infty\}$ مع $x_\infty \equiv 0$
        \emph{(من أجل أي حل، ضع $c =
        \sup\{t : x(t) = 0\}$ وكامل $(\sqrt x)' = 1$
        على $\{x > 0\}$)}، وعلى أن $c \mapsto x_c$ متصل من
        $\intcc0\infty$ (\omterm{def:b3:topology:locallycompact}{بالترصيص بنقطة واحدة}، أي بلصق
        $c = \infty$ بوصفها نهاية) إلى $\mathcal C([0,1])$، واخلص إلى أن
        $\mathcal S$ متراصة ومترابطة فعلًا --- أي قمع على شكل
        قطعة.
  \item (المجموعات القابلة للبلوغ) استنتج من السؤال 20 أنه من
        أجل كل $t$ مثبَّت تكون \emph{المجموعة القابلة للبلوغ}
        $\mathcal S(t) =
        \{x(t) : x \in \mathcal S\}$ متراصة؛ واحسبها من أجل مثال السؤال 21 وتحقق من
        أنها مترابطة أيضًا: $\mathcal S(t) = \intcc0{t^2}$ --- أي إن كل حالة وسيطة
        يبلغها حل ما.
\end{enumerate}

\medskip
\noindent\textbf{الجزء السابع --- تكملات: التعلق، والأمثلية،
ومخطط محسوب يدويًا.}
\begin{enumerate}[resume]
  \item (التعلق الاتصالي) لنفترض أن $f$ ليبشيتزي بثابت $L$ في
        $x$ على $R$، وليكن $x, y$ حلين بقيمتين ابتدائيتين
        $x_0, y_0$ عند $t_0$. بتكييف برهان السؤال 17 مع
        المتراجحة $e(t) \leq \norm{x_0
        - y_0} + \int_{t_0}^tL\,e(s)\dd s$، برهن على
        \[
        \norm{x(t) - y(t)} \leq \norm{x_0 - y_0}\,
        \eu^{L(t - t_0)},
        \]
        وتحقق على $x' = Lx$ من أن الحد مبلوغ: أي إن غرونوال
        دقيقة. واستنتج الوحدانية من جديد ($x_0 =
        y_0$)، وأن تطبيق
        التدفق $x_0 \mapsto x(t; x_0)$ ليبشيتزي بثابت $\eu^{LT}$، حيثما كان
        معرَّفًا.
  \item (أوسغود أمثلية) وبالعكس، لتكن $\omega$ متصلة وغير
        متناقصة وموجبة على $\intoo0\infty$ مع $\omega(r) \to 0$ حين
        $r \to 0^+$ لكن
        \[
        \int_0^1\frac{\dd r}{\omega(r)} < +\infty,
        \]
        ولنمدّد $f(x) = \omega(\abs x)$ بوضع $f(0) = 0$. برهن على أن $\Omega(x) = \int_0^x\frac{\dd r}{\omega(r)}$
        تقابل متزايد من $\intoc0{1}$ على $\intoc0{\Omega(1)}$، وأن مقلوبه
        $g$ يحل $g'
        = \omega(g)$ مع $g(0^+) = 0$، وأن $g$، ممدَّدًا
        بالقيمة $0$ من أجل $t \leq 0$، حلٌّ من الصنف
        $\mathcal C^1$ للمعادلة $x' = f(x)$ مارٌّ بالنقطة
        $(0, 0)$ ومختلف عن $x \equiv
        0$ \emph{(من أجل $g'(0) = 0$،
        حُدّ $\frac{g(t)}t$ بالمقدار $\omega(g(t))$)}. واخلص
        إلى: أن فرضية التباعد في السؤال 15 ليست تسهيلًا بل هي
        الحد الفاصل الدقيق للوحدانية؛ واستعد الجزء الثالث من
        $\omega(r) = 2\sqrt r$، $\Omega(x) = \sqrt x$.
  \item (أويلر محسوبًا يدويًا) من أجل $x' = x$ مع $x(0) = 1$
        على $[0, 1]$: برهن على أن مضلع أويلر يحقق $\varphi_n(1) = \bigl(1 + \frac1n\bigr)^n$.
        وبرهن على النشر
        \[
        \Bigl(1 + \frac1n\Bigr)^{n}
        = \eu\Bigl(1 - \frac1{2n} + O\Bigl(\frac1{n^2}\Bigr)
        \Bigr),
        \]
        ومنه فالخطأ عند $t = 1$ هو $\frac{\eu}{2n} +
        O(n^{-2})$: أي سرعة
        $O(\frac1n)$ من السؤال 18، بالثابت المضبوط. وتحقق
        عدديًا من أجل $n = 10$: $1.1^{10} = 2.59374$ مقابل $\eu \approx 2.71828$، أي بخطأ
        قدره $0.12454$ يُقارن بالمقدار $\frac{\eu}{20}
        \approx 0.13591$.
\end{enumerate}
\end{problem}
